\subsection{The exact scope of a proposed nonzero source pairing}
\label{sec:kt-actual-source-pairing-audit}

This is an audit of a possible next arithmetic argument, not an
additional source exclusion.  Retain the hypotheses and notation
of (CM.1)--(CM.4) and the three auxiliary units.  Put
\(\alpha=\eta\psi\), \(\beta=\psi^{-1}\), and let
\(\mathfrak g\) be the residual trace-zero adjoint module of
the polarized rank-three datum.  The surviving first-order
classes have the form
\(u\in H^1(G_{\mathbf Q,S},k(\alpha))\) and a strict
\(v\in H^1(G_{\mathbf Q,S},k(\beta))\).

\begin{lemma}[KT-SOURCE-TATE-PAIRING-SCOPE]
For every positive integer \(m\),
\[
 \operatorname{Hom}_{G_{\mathbf Q}}
       (\mathfrak g^{\otimes m},k(\omega))=0.
 \tag{TP.1}
\]
In particular, the polarization or an adjoint trace contraction
does not itself provide a Tate-valued bilinear pairing of the
two displayed extension classes.  Their ordinary cup product
has coefficient \(k(\eta)\) and is zero in global cohomology.
\end{lemma}
\begin{proof}
Because \(p\nmid D\), the quadratic character is unramified
at \(p\).  On inertia there, the characters in
\(\mathfrak g\) are powers of \(\omega\) with exponents
\(0,\pm e,\pm2e\).  Choose an inertia element on which
\(\omega\) generates \(\mathbf F_p^\times\).
All these exponents are even.  Every tensor power therefore
has only even exponents, while the target exponent in (TP.1)
is one modulo the even integer \(p-1\).  The inertia
eigenspaces have disjoint characters, proving (TP.1).
This proof uses the actual local residual characters and
does not change upon imposing characteristic-zero
irreducibility on a source point.

Moreover \(\alpha\beta=\eta\), and the checked
quadratic mixed-cup theorem gives
\(H^2(G_{\mathbf Q,S},k(\eta))=0\).
Thus \(u\smile v=0\).  The Cartier dual of
\(k(\beta)\) is \(k(\omega\psi)\), rather than
\(k(\eta\psi)\).  No identification of these
characters is possible on inertia at \(p\).
\end{proof}

The lemma concerns equivariant contractions of the stated
coefficient modules.  It does not rule out a higher
Bockstein or Massey construction with an additional
arithmetic input.

\paragraph{The positive module that actually occurs in duality.}
Let \(U=\operatorname{Spec}\mathbf Z[1/p]\).
The lower relative quintic group of (HQ.13) is
\(H_c^2(U,k(\beta))\), with the usual modified real-place
term, which has no odd-primary contribution.  Artin--Verdier
duality gives a perfect pairing with
\[
 H^1(G_{\mathbf Q,\{p\}},k(\omega\psi)).
 \tag{TP.2}
\]
This assertion uses Milne, \emph{Arithmetic Duality Theorems},
second edition, Theorem II.3.1(b) and Corollary II.3.3,
printed pp.~176--178, together with the compact-support
definition already used in the excision proof.  Those
statements and the full short proof of the corollary were
read for this audit; the global duality theorem is imported.
The trace pairing of \(k/\mathbf F_p\) turns the resulting
finite-group duality into the stated \(k\)-duality.

Writing \(F=\mathbf Q(\mu_p)\), \(\Delta=\operatorname{Gal}(F/\mathbf Q)\),
and \(\mathcal E_F=\mathcal O_F[1/p]^\times\), Kummer
theory gives the exact sequence
\[
\begin{split}
0\longrightarrow&
 (k\otimes_{\mathbf F_p}\mathcal E_F/\mathcal E_F^p)^{(\omega^{-e})}
 \longrightarrow H^1(G_{\mathbf Q,\{p\}},k(\omega\psi))\\
 \longrightarrow&
 (k\otimes_{\mathbf F_p}A_F[p])^{(\omega^{-e})}
 \longrightarrow0.
\end{split}
 \tag{TP.3}
\]
Indeed, restriction to \(F\) and exact \(\Delta\)-invariants
are justified by \(p\nmid p-1\).  Identifying the Kummer
coefficient with \(\mu_p\) twists the character to
\(\omega^{-e}\).  The prime above \(p\) is principal,
so inverting it does not change the class group.
Its contribution to the group of \(p\)-units is a trivial
character and does not affect this nontrivial component.
The Dirichlet logarithmic unit representation contains
each nontrivial even character once.  Its integral
\(\omega^{-e}\)-part is free of rank one; the possible
\(p\)-power root-of-unity torsion has odd character
\(\omega\) and contributes nothing here.
Thus the left term in (TP.3) is a line, and the right
term has the original strict dimension \(d\).
This explains the previously proved dimension \(1+d\),
but supplies no value for the quintic obstruction.

Perfect duality says that a \emph{nonzero} relative
obstruction is detected by some element of (TP.2).
It does not say that irreducibility of a characteristic-zero
point makes a particular obstruction nonzero, nor does it
identify the auxiliary class \(u\) with an element of
the unit line in (TP.3).  The missing statement is an
arithmetic evaluation or nondegeneracy theorem for the
actual higher operation, followed by a verified
source-to-family comparison.

\paragraph{Two exact published nonvanishing tests and their boundaries.}
McCallum--Sharifi, \emph{A cup product in the Galois
cohomology of number fields}, author version dated
5 December 2003, pp.~10--14, explicitly assumes
Vandiver in Sections 4--5.  Lemma 4.1 uses its
cyclotomic-unit index consequence.  Proposition 4.2
expresses pairing with the cyclotomic root-of-unity
direction through the derivative of an odd Iwasawa
characteristic series; the following paragraph obtains
nonzero pairing when that component has lambda
invariant one.  The authors note that cyclicity suffices
for many parts.  The complete proposition proof and
the stated scope were read.  Neither its root-of-unity
input nor its integral hypotheses are supplied by
the auxiliary class \(u\) in our source problem.

Sharifi, \emph{Iwasawa theory and the Eisenstein ideal},
author pp.~6, 11 and 25--27, gives a different genuine
arithmetic test.  Under Hypotheses 1--2, including
\(p\mid B_{1,\theta\omega^{-1}}\) and
\(p\nmid B_{1,\omega\theta^{-1}}\), Theorem 5.6
equates the indicated cup-product surjectivity with
\(U_p-1\) generating the ordinary Eisenstein ideal.
Its entire proof was read, as were the exact
hypotheses and the subsequent finite-range computation.
Theorem 5.7 and Corollary 5.9 establish the required
generation and surjectivity for \(p<1000\).
They do not assert it at every prime.

To use that test on the surviving rank-three source,
one would still need an integral comparison from its
specified off-diagonal module to the indicated
Eisenstein ideal and the relevant generator or
Frobenius value.  The presently checked source
predicates and trace comparisons furnish neither.
A characteristic-zero point or a varying-weight
cuspidal family alone does not show that its
off-diagonal extension class lifts freely across
the reducibility quotient: products of other
off-diagonal entries can occur in its cocycle
equation.  No such cancellation is discarded here.

This last boundary is already visible in the actual matrix
identity.  Restrict a source representation to \(G_E\), put
\(D(g)=\operatorname{diag}(\chi^e(g),1,\chi^{-e}(g))\),
and write \(\rho(g)=(1+A_g)D(g)\) in its integral
split-residual lattice.  Then
\[
 A_{gh}=A_g+D(g)A_hD(g)^{-1}
                    +A_gD(g)A_hD(g)^{-1}.
 \tag{TP.4}
\]
The negative entry \(A_{21}\) therefore satisfies its
linear cocycle equation only after quotienting out the
three terms
\[
 A_{21}(g)A_{11}(h)
 +\chi^{-e}(g)A_{22}(g)A_{21}(h)
 +\chi^{-2e}(g)A_{23}(g)A_{31}(h).
 \tag{TP.5}
\]
They lie in the square of the residual maximal ideal,
but need not vanish in a higher coefficient quotient.
Absolute irreducibility does not set them to zero:
by the characteristic-zero Burnside theorem the linear
span of the matrices \(\rho(g)\) is the full matrix
algebra, so each off-diagonal entry is a nonzero
function of \(g\).  In particular the last product
in (TP.5), with independently varying \(g,h\), is
not identically zero.  This does not assert that
the sum in (TP.5) is a nonzero cohomology class.
It proves why one cannot infer vanishing of the
isolated extension's Bockstein merely from the
existence or irreducibility of the full lift.

This audit therefore proves no additional Selmer
upper bound, source exclusion, or vanishing degree.
